\documentclass[aps,prd,reprint,preprintnumbers,showpacs,floatfix,nofootinbib,superscriptaddress]{revtex4-2}
\usepackage[utf8]{inputenc}
\usepackage{parskip}

\usepackage{amssymb}
\usepackage{amsmath}
\usepackage{mathtools}
\usepackage[dvipsnames]{xcolor}
\usepackage{xspace}
\usepackage{tabularx}
\usepackage{siunitx}
\usepackage{multirow}
\usepackage{subcaption}
\usepackage{graphicx}
\usepackage{xstring}
\usepackage{etoolbox}
\usepackage{notoccite}
\usepackage{natbib}
\usepackage{mathrsfs}
\usepackage{lineno}
\usepackage{tensor}
\usepackage{accents}
\usepackage{microtype}
\usepackage[justification=raggedright,singlelinecheck=false,format=plain]{caption}

\usepackage{tocbasic}
\setcounter{tocdepth}{1}
\DeclareTOCStyleEntry[numwidth=25pt,linefill=\bfseries\TOCLineLeaderFill]{tocline}{section}
\DeclareTOCStyleEntry[entryformat=\textit,numwidth=10pt,linefill=\TOCLineLeaderFill]{tocline}{subsection}
\DeclareTOCStyleEntry[entryformat=\textit,numwidth=10pt,linefill=\TOCLineLeaderFill]{tocline}{subsubsection}

\parindent 2mm

%MacrosJournalNames

% Journal names
\def\aj{\rm{AJ}}
\def\actaa{\rm{Acta Astron.}}
\def\araa{\rm{ARA\&A}}
\def\apj{\rm{ApJ}}
\def\apjl{\rm{ApJ}}
\def\apjs{\rm{ApJS}}
\def\ao{\rm{Appl.~Opt.}}
\def\apss{\rm{Ap\&SS}}
\def\aap{\rm{A\&A}}
\def\aapr{\rm{A\&A~Rev.}}
\def\aaps{\rm{A\&AS}}
\def\azh{\rm{AZh}}
\def\baas{\rm{BAAS}}
\def\bac{\rm{Bull. astr. Inst. Czechosl.}}
\def\caa{\rm{Chinese Astron. Astrophys.}}
\def\cjaa{\rm{Chinese J. Astron. Astrophys.}}
\def\icarus{\rm{Icarus}}
\def\jcap{\rm{J. Cosmology Astropart. Phys.}}
\def\jrasc{\rm{JRASC}}
\def\memras{\rm{MmRAS}}
\def\mnras{\rm{MNRAS}}
\def\na{\rm{New A}}
\def\nar{\rm{New A Rev.}}
\def\pra{\rm{Phys.~Rev.~A}}
\def\prb{\rm{Phys.~Rev.~B}}
\def\prc{\rm{Phys.~Rev.~C}}
\def\prd{\rm{Phys.~Rev.~D}}
\def\pre{\rm{Phys.~Rev.~E}}
\def\prl{\rm{Phys.~Rev.~Lett.}}
\def\pasa{\rm{PASA}}
\def\pasp{\rm{PASP}}
\def\pasj{\rm{PASJ}}
\def\rmxaa{\rm{Rev. Mexicana Astron. Astrofis.}}
\def\qjras{\rm{QJRAS}}
\def\skytel{\rm{S\&T}}
\def\solphys{\rm{Sol.~Phys.}}
\def\sovast{\rm{Soviet~Ast.}}
\def\ssr{\rm{Space~Sci.~Rev.}}
\def\zap{\rm{ZAp}}
\def\nat{\rm{Nature}}
\def\iaucirc{\rm{IAU~Circ.}}
\def\aplett{\rm{Astrophys.~Lett.}}
\def\astropj{\rm{Astrophys.~J.}}
\def\apspr{\rm{Astrophys.~Space~Phys.~Res.}}
\def\bain{\rm{Bull.~Astron.~Inst.~Netherlands}}
\def\fcp{\rm{Fund.~Cosmic~Phys.}}
\def\gca{\rm{Geochim.~Cosmochim.~Acta}}
\def\grl{\rm{Geophys.~Res.~Lett.}}
\def\jcp{\rm{J.~Chem.~Phys.}}
\def\jgr{\rm{J.~Geophys.~Res.}}
\def\jqsrt{\rm{J.~Quant.~Spec.~Radiat.~Transf.}}
\def\memsai{\rm{Mem.~Soc.~Astron.~Italiana}}
\def\nphysa{\rm{Nucl.~Phys.~A}}
\def\nphys{\rm{Nucl.~Phys.}}
\def\physrep{\rm{Phys.~Rep.}}
\def\physscr{\rm{Phys.~Scr}}
\def\planss{\rm{Planet.~Space~Sci.}}
\def\procspie{\rm{Proc.~SPIE}}
\def\procrsa{\rm{Proc.~Math.~Phys.~Eng.~Sci.}}
\def\annmath{\rm{Ann.~Math.}}
\def\mathann{\rm{Math.~Ann.}}
\def\annphys{\rm{Ann.~Physics (N.~Y.)}}
\def\annderphys{\rm{Ann.~Phys.}}
\def\ptrsla{\rm{Philos.~Trans.~R.~Soc~A}}
\def\ptrslo{\rm{Philos.~Trans.~R.~Soc~I}}
\def\jmp{\rm{J.~Math.~Phys.}}
\def\cqg{\rm{Class.~Quantum~Gravity}}
\def\physrev{\rm{Phys.~Rev.}}
\def\foundphys{\rm{Found.~Phys.}}
\def\epjc{\rm{Eur.~Phys.~J.~C}}
\def\natast{\rm{Nat.~Astron.}}
\def\pdu{\rm{Phys.~Dark~Universe}}
\def\araa{\rm{Annu.~Rev.~Astron.~Astrophys.}}
\def\nrp{\rm{Nat.~Rev.~Phys.}}
\def\rpp{\rm{Rep.~Prog.~Phys.}}
\def\ijmpd{\rm{Int.~J.~Mod.~Phys.~D}}
\def\pletta{\rm{Phys.~Lett.~A}}
\def\plettb{\rm{Phys.~Lett.~B}}
\def\grg{\rm{Gen.~Relativ.~Gravit.}}
\def\phyrep{\rm{Phys.~Rep.}}
\def\jhep{\rm{J.~High~Energy~Phys.}}
\def\rmphy{\rm{Rev.~Mod.~Phys.~D}}
\def\rmphyo{\rm{Rev.~Mod.~Phys}}
\def\fass{\rm{Front.~Astron.~Space~Sci.}}
\def\progtp{\rm{Prog.~Theor.~Phys.}}
\def\appb{\rm{Acta~Phys.~Pol.~B}}
\def\crphys{\rm{C.~R.~Phys.}}
\def\ijtp{\rm{Int.~J.~Theor.~Phys.}}
\def\progpnp{\rm{Prog.~Part.~Nucl.~Phys.}}
\def\nucphyb{\rm{Nucl.~Phys.~B}}
\def\npcs{\rm{Nonlin.~Phenom.~Complex~Syst.}}
\def\nar{\rm{New.~Astron.~Rev.}}
\def\cjp{\rm{Chin.~J.~Phys.}}
\def\ujp{\rm{Ukr.~J.~Phys.}}
\def\aaca{\rm{Adv.~Appl.~Clifford~Algebras}}
\def\sjetpl{\rm{Sov.~JETP~Lett.}}
\def\mathz{\rm{Math.~Z.}}
\def\macsp{\rm{Mem.~Acad.~St.~Petersbourg}}
\def\rprogphy{\rm{Rept.~Prog.~Phys.}}
\def\annhp{\rm{Ann.~Henri~Poincar\'e}}
\def\astronj{\rm{Astron.~J.}}
\def\ijgmmp{\rm{Int.~J.~Geom.~Meth.~Mod.~Phys.}}
\def\pispf{\rm{Proc.~Int.~Sch.~Phys.~Fermi}}
\def\livrevrel{\rm{Living~Rev.~Relativ.}}


\usepackage{hyperref}
\hypersetup{%
 colorlinks = true,%
 linkcolor = red,%
 citecolor = red,%
 filecolor = red,%
 urlcolor = red%
 }%
\usepackage[capitalize]{cleveref}
\AddToHook{cmd/appendix/before}{\crefalias{section}{appendix}}
\makeatletter
\renewcommand{\paragraph}{%
\@startsection{paragraph}{4}%
{\z@}{1.21ex \@plus 1ex \@minus .2ex}{0.9em}%
{\normalfont\normalsize\bfseries}%
}
\usepackage{titlesec}
\newrobustcmd{\pea}[1]{%
	\emph{#1}\textbf{\ \ \ ---}
}
\titleformat{\paragraph}[runin]{\normalfont\normalsize\bfseries}{\emph\theparagraph}{1em}{\pea}
\titleformat{\section}[block]{\normalfont\bfseries\centering}{\MakeUppercase\thesection}{1em}{\MakeUppercase}
\makeatother

\newcommand{\hunit}{km\,s$^{-1}$\,Mpc$^{-1}\,~$}
\newcommand{\LCDM}{\ensuremath{\Lambda\mathrm{CDM}}}
\newcommand{\Mp}{M_\mathrm{P}^2}
\newcommand{\RR}{\mathcal{R}}
\newcommand{\TT}{\mathcal{T}}
\newcommand{\PGTq}{\ensuremath{\text{PGT}^{\text{q}+}}}
\DeclareSIUnit\Mpc{Mpc}

% Comment macros for authors
\newcommand{\WB}[1]{\textbf{\textit{\textcolor{blue}{WB: #1}}}}
\newcommand{\SL}[1]{\textbf{\textit{\textcolor{red}{SL: #1}}}}
\newcommand{\WH}[1]{\textbf{\textit{\textcolor{ForestGreen}{WH: #1}}}}

\newrobustcmd{\SmashAcute}[1]{%
	{\vphantom{#1}\smash{\Acute{#1}}}
}
\newrobustcmd{\ECT}[1]{%
	\tensor{\mathcal{T}}{#1}
}
\newrobustcmd{\ECR}[1]{%
	\tensor{\mathcal{R}}{#1}
}
\newrobustcmd{\MetricPerturbation}[1]{%
	\tensor*{h}{#1}
}
\newrobustcmd{\RiemannianR}[1]{%
	\tensor{R}{#1}
}
\newrobustcmd{\StressEnergyTensor}[1]{%
	\tensor{\mathbb{T}}{#1}
}
\newrobustcmd{\TorsionSource}[1]{%
	\tensor{\Delta}{#1}
}
\newrobustcmd{\AffineConnection}[1]{%
	\tensor{\Gamma}{#1}
}
\newrobustcmd{\LeviCivitaConnection}[1]{%
	\big\{\tensor{}{#1}\big\}
}
\newrobustcmd{\PD}[1]{%
	\tensor{\partial}{#1}
}
\newrobustcmd{\CD}[1]{%
	\tensor{\nabla}{#1}
}
\makeatletter

\allowdisplaybreaks

\begin{document}

\title{Investigating the cosmological properties of the Stigma Q-torsion model}

\author{Sinah Legner}
\email{sl2091@cam.ac.uk}
\affiliation{Astrophysics Group, Cavendish Laboratory, JJ Thomson Avenue, Cambridge CB3 0HE, UK}
\affiliation{Kavli Institute for Cosmology, Madingley Road, Cambridge CB3 0HA, UK}

\author{Will Barker}
\email{barker@fzu.cz}
\affiliation{Central European Institute for Cosmology and Fundamental Physics, Institute of Physics of the Czech Academy of Sciences, Na Slovance 1999/2, 182 00 Prague 8, Czechia}

\author{Will Handley}
\email{wh260@cam.ac.uk}
\affiliation{Kavli Institute for Cosmology, Madingley Road, Cambridge CB3 0HA, UK}
\affiliation{Institute of Astronomy, Madingley Road, Cambridge CB3 0HA, UK}

\begin{abstract}
	Abstract text goes here.
\end{abstract}

\maketitle

% Reduce spacing in table of contents
{
\setlength{\parskip}{0pt}
\setlength{\topsep}{0pt}
\setlength{\partopsep}{0pt}
\setlength{\itemsep}{0pt}
\renewcommand{\baselinestretch}{0.8}
\tableofcontents
}
\setlength{\parskip}{6pt}

\section{Introduction}\label{sec:intro}

\paragraph*{Poincar\'e Gauge Theory} The Poincar\'e Gauge Theory (PGT) of gravity provides a natural framework for introducing torsion and curvature as independent dynamical quantities in spacetime. Within the catalogue of ghost-free, tachyon-free, and power-counting renormalisable PGT models identified in~\cite{Lin:2018awc, Lin:2019ugq}, the \emph{Stigma} model represents a distinct branch characterised by a non-zero Q-torsion sector. The general PGT Lagrangian admits a rich parameter space, of which the Stigma model occupies a specific region defined by the coupling constants $(\varsigma, \upsilon_2)$.

\begin{equation}
		\begin{align}
				\mathcal{L} = & \frac{3 \Mp \upsilon_2}{4 \varsigma}\Big[ \frac{1}{2} \big( 4 \varsigma - \varpi^2\big) R - g^{\mu\nu}\partial_\mu \varpi \partial_\nu \varpi \\
							  & + \varsigma \phi^2 - \varpi^2 \phi^2 + \frac{4}{3}\varsigm \sqrt{J_\mu J^\mu} \Big] - \Mp \Lambda + \mathcal{L}_\mathrm{M}\\
						J_\mu \equiv & \varpi^3 \partial_\mu (\phi/\varpi) - \partial_\mu \phi
		\end{align}
\end{equation}

\begin{equation}
		\begin{align}
				H^2 = & \frac{1}{\varpi^2-4\varsigma} \\
					  & \times \Bigg[-\frac{4 H_0^2\varsigma}{\upsilon_2}\Big(\Omega_{\mathrm{r},0}a^{-4}+\Omega_{\mathrm{m},0}a^{-3}+\Omega_{\Lambda,0}\Big) \\
						& + 2 H \varpi\dot{\varpi} + \dot{\varpi}^2  + \frac{4\left(H\varsigma\left(\varpi^2-1\right)+\varsigma\varpi\dot{\varpi}\right)^2}{\varsigma-\varpi^2}\Bigg]
						\end{align}
\end{equation}

\begin{equation}
		\begin{aligned}
				\ddot{\varpi} & =\frac{1}{a^2\left(\varsigma-\varpi^2\right)\left(\varsigma+\left(4\varsigma^2-1\right)\varpi^2\right)} \\
							  & \times \Bigg[-3\varsigma^2 a \dot a \dot{\varpi} + 6\varsigma\left(1-2\varsigma^2\right)a\varpi^2\dot a\dot{\varpi} \\
							  & +3\left(4\varsigma^2-1\right)a\varpi^4\dot a\dot{\varpi}+\left(4\varsigma^2-1\right)\varpi^5\left(\dot a^2+a\ddot a\right) \\
							  &+2\varsigma\varpi^3\Big(\left(1-4\varsigma\right)\dot a^2+\left(1-2\varsigma(1+\varsigma)\right)a\ddot a\Big)\\
							  & +\varsigma^2\varpi\Big(3\dot a^2+a\left(-4\varsigma a\dot{\varpi}^2+\left(4\varsigma-1\right)\ddot a\right)\Big) \Bigg]
				
		\end{aligned}
\end{equation}





\section{Cosmology}\label{sec:cosmology}

\subsection{Behavior at radiation-dominated epoch}
We look at the behaviour of the field equations at the radiation dominated epoch, by looking only at a fluid-dominated epoch, i.e. $H = 2/(3(1+w_i)t)$, where $w_i$ is the equation of state parameter for the fluid. We find that from the field equation of $\varpi$, which we assume the field $\varpi$ is a constant. 
\begin{equation}
		\label{eq:varpi_field_eq}
		\begin{align}
				\left.\frac{\partial L}{\partial \varpi}\right|_{\varpi=\varpi_i} \propto & \quad \frac{\varpi_i}{(w_i + 1)^2 (\varpi_i^2 - \varsigma)^2 \varsigma}\\
																								   &\times \Biggl[(3 w_i-1)(4 \varsigma^2 - 1) \varpi_i^4 \\
																								& + \varsigma^2\Big( -7 + 4 \varsigma +3 w_i (4 \varsigma -1) \\
																								& -2 \varpi_i^2 \varsigma \big(1+ 2(\varsigma - 3) \varsigma \\
																								& + w_i (6\varsigma (1+\varsigma)-3)\big)\Big)\Biggr]
		\end{align}
\end{equation}
Solving $\varpi^2$ for this equation reveals that
\begin{equation}
		\begin{align}
				\varpi_i^2 = & \frac{\varsigma}{(3 w_i - 1)(4 \varsigma^2 - 1)} \\
						   & \times \Big[1+2(\varsigma-3)\varsigma + w_i \Big(6\varsigma (1+\varsigma)-3\Big) \\
						   & \pm 2 \sqrt{(\varsigma-1)^2 \Big(2 - 6 w_i + (\varsigma + 3 w_i \varsigma)^2\Big)}\Big] 
		\end{align}
\end{equation}
We can see that the solution contain poles at $w_i = 1/3$ and $\varsigma = 1/2$. However, upon further investigation we realised that the pole of $w_i = 1/3$ is a false-pole. Subsituting $w_i = 1/3$ at~\cref{eq:varpi_field_eq} we find that the equation reduces to
\begin{equation}
		\left.\frac{\partial L}{\partial \varpi}\right|_{\varpi=\varpi_i}^{w_i=1/3} \propto -\frac{(\varpi_i^2 - 1)(\varsigma-1)\varpi_i \varsigma }{(\varpi_i^2 - \varsigma)^2}
\end{equation}
The numerator shows that at radiation-dominated epoch, $\varpi_i = 1$ and $\varsigma = 1$ are the only two solutions. To ensure that the field equation goes to zero, one of them must be set to be true.

\subsection{Power series}
\begin{equation}
		\begin{align}
				\varpi = & 1 \mp \frac{H_0 \Omega_{\mathrm{m, 0}}\sqrt{4 \varsigma - 1}}{8 \sqrt{- \upsilon_2 \Omega_{\mathrm{r, 0}} \varsigma}} \tau\\
						 & - \frac{3 H_0^2 \Omega_{\mathrm{m, 0}}^2 (4 \varsigma - 1)(2 + \varsigma(3 \varsigma - 2))}{128 \upsilon_2 \Omega_{\mathrm{r, 0}}(\varsigma - 1)\varsigma (8 \varsigma^2 + \varsigma -3)}\tau^2\\
						 & \pm \frac{H_0^3\Omega_\mathrm{m, 0}^3}{1024(\varsigma-1)^2}\Big(\frac{4 \varsigma -1}{- \upsilon_2 \varsigma \Omega_\mathrm{r,0}}\Big)^{3/2} \\ 
						 & \times \frac{\tau^3}{(8 \varsigma^2 + \varsigma -3)(20 \varsigma^2 + \varsigma -6)} \times \Big(192 \varsigma^6 \\
						 &+ 380 \varsigma^5 -720 \varsigma^4 -225 \varsigma^3 +175 \varsigma^2 + 72 \varsigma - 36\Big)\\
						 &+ \mathcal{O}(\tau^4)
		\end{align}
\end{equation}
\begin{equation}
		\begin{align}
				a = & \pm \frac{2 H_0 \sqrt{\Omega_{\mathrm{r,0}} \varsigma}}{\alpha \sqrt{\upsilon_2 (1- 4 \varsigma)}}\tau - \frac{H_0^2 \Omega_\mathrm{m, 0}}{4 \alpha \upsilon_2} \tau^2\\
					& \pm \frac{3 H_0^3 \Omega_{\mathrm{m, 0}}^2 \sqrt{\varsigma (4 \varsigma - 1)}(\varsigma(24 \varsigma + 5) -2)}{64 \alpha (-\upsilon_2)^{3/2} \sqrt{\Omega_{\mathrm{r, 0}}}(\varsigma - 1)(8 \varsigma^2 + \varsigma -3)}\tau^3\\
					& + \frac{H_0^4 \Omega_\mathrm{m,0}^3 (4 \varsigma - 1)}{512 \alpha \upsilon_2^2 \Omega_\mathrm{r,0} (\varsigma - 1)^2}\\
					& \times \frac{240\varsigma^5-1510\varsigma^4-246\varsigma^3+357\varsigma^2+31\varsigma-6}{(8 \varsigma^2 + \varsigma -3)(20 \varsigma^2 + \varsigma -6)}\tau^4\\
					& + \mathcal{O}(\tau^5)
		\end{align}
\end{equation}

\begin{equation}
				  \begin{align}
				  H^2 ={}& \frac{1}{4\varsigma-\varpi^2} \Bigg[ \frac{4H_0^2\varsigma}{-\upsilon_2^2} \left( \Omega_{\mathrm r,0}a^{-4} +\Omega_{\mathrm m,0}a^{-3} +\Omega_{\Lambda,0} \right) + \\
						 &2H\varpi\dot{\varpi} +\dot{\varpi}^{\,2} \nonumber\\ 
						 &+ \frac{ 4\left[ H\varsigma\left(\varpi^2-1\right) +\varsigma\varpi\dot{\varpi} \right]^2 }{ \varsigma-\varpi^2 } \Bigg]. \label{eq:TorCF1} \end{align}
				 \end{equation}

\section{Methodology}\label{sec:methodology}

Blahblahblahblahblah

\section{blah}

Methodology text goes here.

\SL{Describe the numerical ODE solver, initial conditions from power-series expansion, rescaling of the scale factor, and comparison with \LCDM{}.}

\section{Results}\label{sec:results}

Results text goes here.

\SL{Present plots of the expansion history, equation of state, and comparison with \LCDM{}.}

\section{Discussion}\label{sec:discussion}

Discussion text goes here.

\section{Conclusion}\label{sec:conclusion}

Conclusion text goes here.

\begin{acknowledgments}
Acknowledgments go here.
\end{acknowledgments}

\bibliographystyle{apsrev4-2}
\bibliography{paper_Stigma}

\end{document}
